
\documentclass[11pt]{article}
\usepackage[utf8]{inputenc}
\usepackage{amsmath,amssymb}
\usepackage{enumitem}
\usepackage{geometry}
\geometry{margin=1in}

% Custom macros for questions
\newcommand{\QBlock}[3]{%
  \subsection*{Question #1}
  #2
  \begin{enumerate}[label=(\Alph*)]
  #3
  \end{enumerate}
  \vspace{1em}
}

\newcommand{\Correct}[2]{%
  \item \textbf{#1} \hfill $\checkmark$ \\
  \textit{Explanation: #2}
}

\newcommand{\Wrong}[2]{%
  \item #1 \\
  \textit{Explanation: #2}
}

\begin{document}

\title{Isomorphic Physics Questions}
\author{Generated by AI}
\date{\today}
\maketitle

% Questions go here using \QBlock, \Correct, \Wrong macros

\QBlock{1}{Which of the following assumptions is valid in physical analysis?}{
  \Correct{We can assume a "point mass" when analyzing a car's motion over a long highway.}{For large-scale motion, treating a car as a point mass simplifies the analysis.}
  \Wrong{We can assume a "cubical car" when analyzing the turning radius of the car.}{A cubical assumption is unrealistic for this physical situation.}
  \Wrong{Ignoring air resistance means assuming the car is weightless.}{Ignoring air resistance doesn't imply weightlessness.}
  \Wrong{None of the above is correct.}{Option A is actually a correct simplification.}
}

\QBlock{2}{Which of the following conversions is correct?}{
  \Wrong{1 penny = 1 dollar}{A penny is not equal to a dollar.}
  \Correct{10 cents = 0.1 dollars}{Ten cents is indeed 0.1 dollars as $0.1 \$/\text{cent(s)}.$}
  \Wrong{1 = 0.1}{One is not equal to 0.1.}
  \Wrong{100 pennies = 0.1 dollars}{100 pennies equal 1.00 dollars, not 0.1.}
}

\QBlock{3}{Which option corresponds to an incorrect unit for the given quantity?}{
  \Wrong{Volume: liter}{Liter is the correct unit for volume.}
  \Correct{Speed: meter/hour}{Speed should be measured in meters per second (m/s), not meters per hour.}
  \Wrong{Pressure: pascal}{Pascal is the correct unit for pressure.}
  \Wrong{Power: watt}{Watt is the correct unit for power.}
}

\QBlock{4}{A ship sets sail from a port at coordinates (0,0) to an island located at $\mathbf{r} = 4\mathbf{i} + 7\mathbf{j}$ km. Which of the following statements is correct?}{
  \Correct{The distance from the port to the island is 8 km.}{Using the Pythagorean theorem, the distance is calculated as $\sqrt{4^2 + 7^2} \approx 8$ km.}
  \Wrong{The distance from the port to the island is 7 km.}{Incorrect calculation based simply on one axis.}
  \Wrong{The distance from the port to the island is 4 km.}{Incorrect value taken from one component.}
  \Wrong{None of the above is correct.}{Option A is correct.}
}

\QBlock{5}{If the displacement of an object is given by the function $\vec{s} = ((2.0\, \text{m/s}^3)t^3 + (1.5\, \text{m/s}^2)t^2)\mathbf{i}$, which of the following is correct?}{
  \Wrong{This describes a motion with constant acceleration.}{The acceleration depends on time, indicating it is not constant.}
  \Wrong{The magnitude of the acceleration is 1.5 m/s$^2$}{The acceleration is not constant; it varies with time.}
  \Correct{The magnitude of the acceleration at $t = 0$ is 0 m/s$^2$}{At $t = 0$, the acceleration is found to be 0 by differentiating the velocity.}
  \Wrong{The magnitude of the acceleration at $t = 0$ is 3.0 m/s$^2$}{The acceleration at $t = 0$ is not 3.0 m/s$^2$.}
}

\QBlock{6}{An object takes 4 seconds to move from position $\vec{r}_A = -3.0\mathbf{i} + 8\mathbf{j}$ to position $\vec{r}_B = 3.0\mathbf{i} + 3.0\mathbf{j}$, where length units are in meters. Which statement is correct?}{
  \Correct{The moving distance is 9.2 m.}{Displacement is calculated using $\sqrt{(3 - (-3))^2 + (3 - 8)^2} \approx 9.2$ m.}
  \Wrong{The speed at $t = 0$ must be 2.3 m/s.}{Speed calculation is inaccurate.}
  \Wrong{The acceleration at $t = 0$ must be 1.05 m/s$^2$}{Insufficient information to determine acceleration.}
  \Wrong{None of the above is correct.}{Option A is correct.}
}

\QBlock{7}{A football is kicked with a speed $v_0$ at an angle $45^\circ$ above the horizontal. What is the speed of the ball at its maximum height?}{
  \Wrong{$v_0$}{At maximum height, the vertical speed is zero, leaving only the horizontal component.}
  \Correct{$v_0 \cos 45^\circ$}{At maximum height, the speed equals the horizontal component $v_0 \cos 45^\circ$.}
  \Wrong{$v_0 \sin 45^\circ$}{The vertical component becomes zero at maximum height.}
  \Wrong{0}{Only the vertical component is zero at the maximum height.}
}

\QBlock{8}{A drone starts from rest and accelerates with an acceleration $\vec{a} = (3.0 \text{ m/s}^2)\mathbf{i} + (1.5 \text{ m/s}^2)\mathbf{j}$. Which statement is correct?}{
  \Wrong{The x-component of velocity remains the same during acceleration.}{The x-component increases due to the positive acceleration.}
  \Wrong{The y-component of velocity remains constant.}{The y-component increases due to the exerted acceleration.}
  \Correct{Both direction and speed change during the motion.}{Both direction and magnitude of velocity change due to provided acceleration.}
  \Wrong{None of the above is correct.}{Option C is actually correct.}
}

\QBlock{9}{If the magnitudes of both masses are tripled and their separation remains constant, what happens to the gravitational force between them?}{
  \Wrong{It stays the same.}{Tripling masses increases the force.}
  \Correct{It increases by a factor of 9.}{Force is proportional to the product of the masses (gravitational force $F \propto m_1 m_2$).}
  \Wrong{It decreases by a factor of $3^{1/2}$.}{The force actually increases.}
  \Wrong{It doubles.}{The factor should be related to squared change in mass.}
  \Wrong{It remains constant while masses increase.}{Force must increase with increased masses.}
}

\QBlock{10}{A launching platform uses a catapult with spring constant $k$ and stretch $x$ to launch projectiles. Which of the following affects the escape speed?}{
  \Wrong{Increasing the mass of the projectile decreases escape speed directly.}{Escape speed doesn't directly depend on mass of the escape object.}
  \Correct{Decreasing the tension in the spring decreases the escape speed.}{Decreasing spring energy (proportional to $kx^2$) lowers the kinetic energy transfer, reducing escape speed.}
  \Wrong{Changing projectile surface area doesn't affect escape speed.}{Surface area isn't relevant in spring bound energy.}
  \Wrong{Different pathways with same spring force result in different speeds.}{Only spring constant and deformation are relevant.}
}

\QBlock{11}{A force $\mathbf{F} = (5.00\, \text{N})\mathbf{i} + (3.00\, \text{N})\mathbf{j}$ acts on a particle moving along the $x$-axis from $x = 1.50$ m to $x = 6.50$ m. What is the work done by the force on the particle in joules?}{
  \Correct{20.0 J}{Work done is the integral of the $x$-component force over the displacement, giving $5 \text{ N} \times (6.5 - 1.5) \text{ m} = 25 \text{ J}$.}
  \Wrong{15.0 J}{Incorrect calculation based on partial force component.}
  \Wrong{0 J}{Non-zero force results in non-zero work.}
  \Wrong{35.0 J}{Work is lower due to $x$ displacement integration only.}
}

\QBlock{12}{A 5.5 kg cart initially moving at 3.1 m/s collides with a wall and rebounds at a speed of 2.0 m/s in backward motion. The collision lasts for 0.25 seconds. What is the average force exerted on the cart during the collision?}{
  \Correct{75.6 N.}{The average force is calculated using $F = \frac{\Delta p}{\Delta t} = \frac{m(v_f - v_i)}{\Delta t}$.}
  \Wrong{20.0 N.}{Calculation based on incorrect change in velocity.}
  \Wrong{100.0 N.}{Force calculation exceeds due to overestimation of change in velocity.}
  \Wrong{None of the above is correct.}{Option A provides correct average force.}
}

\QBlock{13}{An 80 kg basketball player catches a 0.2 kg ball moving at 25 m/s. What is the player's speed after catching the ball?}{
  \Correct{0.0625 m/s}{Using conservation of momentum, $mv = (m_{player} + m_{ball})v_f$, solve for $v_f$.}
  \Wrong{0.100 m/s}{Calculation error in defining the proportionality correctly.}
  \Wrong{125 m/s}{Irregular calculation leading to overestimation.}
  \Wrong{0.114 m/s}{Incorrect balancing of momentum.}
}

\QBlock{14}{A 2,000 kg car traveling at 60 km/h comes to a stop over 25 meters. What is the average braking force applied?}{
  \Correct{9,600 N}{Using $F = ma$ and $v^2 = u^2 + 2as$ to find deceleration.}
  \Wrong{12,800 N}{Overestimation caused by incorrect substitution in the equation.}
  \Wrong{7,500 N}{Underestimation from missing conversion factor.}
  \Wrong{5,000 N}{Calculation error in solving force.}
}

\QBlock{15}{A marble is released from a height of 100 meters. Ignoring air resistance, what is the speed of the marble just before hitting the ground?}{
  \Correct{44.3 m/s}{Using $v = \sqrt{2gh}$, where $g = 9.8 \text{ m/s}^2$ and $h = 100 \text{ m}$.}
  \Wrong{50.0 m/s}{Overestimation in free fall calculation.}
  \Wrong{40.0 m/s}{Underestimation due to incorrect $gh$ substitution.}
  \Wrong{47.5 m/s}{Incorrect calculation exceeding accurate terminal speed.}
}

\QBlock{16}{A constant force $\vec{F} = (4.00\, \text{N})\mathbf{i} - (2.00\, \text{N})\mathbf{j}$ acts on an object moving from $x = 2.0$ m to $x = 8.0$ m. What is the work done by this force?}{
  \Correct{24.0 J}{Work equals the dot product of force along direction $W = 4 \times (8 - 2) = 24$ J.}
  \Wrong{32.0 J}{Calculation based on incorrect $y$ component influence.}
  \Wrong{12.0 J}{Underestimation from main linear force component.}
  \Wrong{0 J}{Non-zero force and displacement give non-zero work.}
}

\QBlock{17}{A 5.0 kg kart initially moving at 4.0 m/s collides into a wall and rebounds at a speed of 1.5 m/s in backward direction. The collision time is 0.1 seconds. What is the average force on the kart during the collision?}{
  \Correct{300 N.}{Applying $F = \frac{\Delta p}{\Delta t} = \frac{m(v_f - v_i)}{\Delta t}$ to find force.}
  \Wrong{200 N.}{Error in calculated force equation.}
  \Wrong{500 N.}{Exaggeration based on velocity changes overly.}
  \Wrong{Zero force operation.}{Force necessary to achieve velocity change.}
}

\QBlock{18}{A 75 kg soccer goalkeeper stops a ball of 0.15 kg moving at 30 m/s. What is the resultant speed of the goalkeeper after the save?}{
  \Correct{0.060 m/s}{Using momentum conservation $mv_{ball} = (m_{goalkeeper} + m_{ball})v_f$.}
  \Wrong{0.045 m/s}{Underestimation due to incorrect velocity assumption.}
  \Wrong{20 m/s}{Unreal overcalculation exceeding realistic speed.}
  \Wrong{0.050 m/s}{Incorrect reduction constant used in calculations.}
}

\end{document}
